\documentclass[12pt,a4paper]{article}

\usepackage[margin=1in]{geometry}
\usepackage{amssymb}
\usepackage{booktabs}
\usepackage{hyperref}
\usepackage{enumitem}
\usepackage{titlesec}
\usepackage{fancyhdr}
\usepackage{xcolor}
\usepackage{array}
\usepackage{tabularx}
\usepackage{multirow}
\usepackage{amsmath}

% ── Colours ───────────────────────────────────────────────
\definecolor{headblue}{HTML}{1A5276}

% ── Section styling ───────────────────────────────────────
\titleformat{\section}
  {\Large\bfseries\color{headblue}}
  {\thesection.}{0.5em}{}[\vspace{2pt}\titlerule]
\titleformat{\subsection}
  {\large\bfseries\color{headblue!80}}
  {\thesubsection}{0.5em}{}

% ── Header / Footer ──────────────────────────────────────
\pagestyle{fancy}
\setlength{\headheight}{14.5pt}
\addtolength{\topmargin}{-2.5pt}
\fancyhf{}
\lhead{\small\textit{Cost Estimation --- Traffic Prediction System}}
\rhead{\small\textit{Group E-16}}
\cfoot{\thepage}

\begin{document}

% ═══════════════════════════════════════════════════════════
% TITLE PAGE
% ═══════════════════════════════════════════════════════════
\begin{titlepage}
\centering
\vspace*{2cm}
{\Huge\bfseries\color{headblue} Software Cost Estimation\\[12pt]}
{\LARGE Function Point Analysis \& Use Case Points\\[8pt]}
{\large Traffic Prediction and Analysis System\\[20pt]}
\rule{0.6\textwidth}{1pt}\\[20pt]

{\large\textbf{Assignment No. 8}}\\[30pt]

\begin{tabular}{ll}
\textbf{Group:} & E-16 \\[4pt]
\textbf{Members:} & Harshal Jaiswal \\
 & Pranav Jagtap \\
 & Suyash Jobanputra \\
 & Shreyash Jobanputra \\
\end{tabular}

\vfill
{\small Department of Computer Engineering\\
\today}
\end{titlepage}

% ═══════════════════════════════════════════════════════════
\section{Problem Statement}
% ═══════════════════════════════════════════════════════════
Perform software cost estimation exercises using \textbf{Function Point Analysis (FPA)} and \textbf{Use Case Points (UCP)} for the Traffic Prediction and Analysis System. Estimate the software size, effort, and approximate development cost based on these two industry-standard techniques.

% ═══════════════════════════════════════════════════════════
\section{Introduction}
% ═══════════════════════════════════════════════════════════
Software cost estimation is a fundamental activity in project management that predicts the resources, time, and budget required to develop a software system. Accurate estimation helps prevent budget overruns, enables proper resource allocation, and sets realistic delivery expectations.

Two widely used estimation techniques are:
\begin{enumerate}[noitemsep, label=\roman*)]
    \item \textbf{Function Point Analysis (FPA)} --- measures software size based on the functional requirements delivered to the user, independent of the programming language or technology used.
    \item \textbf{Use Case Points (UCP)} --- estimates effort based on the complexity of use cases and the environmental and technical factors of the project.
\end{enumerate}

This assignment applies both techniques to the \textbf{Traffic Prediction and Analysis System}, a full-stack web application built using React.js, Node.js/Express, and MongoDB that integrates external APIs and a MobileNet ML model for traffic forecasting and pothole detection.

% ═══════════════════════════════════════════════════════════
% PART A: FUNCTION POINT ANALYSIS
% ═══════════════════════════════════════════════════════════
\section{Function Point Analysis (FPA)}

\subsection{Overview}
Function Point Analysis, introduced by Allan Albrecht at IBM in 1979, quantifies software functionality from the user's perspective. It classifies system interactions into five categories:
\begin{itemize}[noitemsep]
    \item \textbf{External Inputs (EI)} --- data entering the system from outside.
    \item \textbf{External Outputs (EO)} --- data or reports leaving the system.
    \item \textbf{External Inquiries (EQ)} --- input/output combinations for data retrieval.
    \item \textbf{Internal Logical Files (ILF)} --- data maintained within the system.
    \item \textbf{External Interface Files (EIF)} --- data referenced from external systems.
\end{itemize}

\subsection{Step 1: Identify and Classify Function Points}

\subsubsection*{External Inputs (EI)}
\begin{center}
\begin{tabular}{clc}
\toprule
\textbf{\#} & \textbf{Function} & \textbf{Complexity} \\
\midrule
1 & User Registration & Low \\
2 & User Login (JWT Authentication) & Low \\
3 & Submit Traffic Prediction Request & Average \\
4 & Upload Pothole Image & High \\
5 & Submit Infrastructure Complaint & Average \\
6 & Register Event Space & Low \\
7 & Add Construction Data (Admin) & Average \\
8 & Add Traffic Diversion (Admin) & Average \\
9 & Update User Role (Admin) & Low \\
\bottomrule
\end{tabular}
\end{center}

\subsubsection*{External Outputs (EO)}
\begin{center}
\begin{tabular}{clc}
\toprule
\textbf{\#} & \textbf{Function} & \textbf{Complexity} \\
\midrule
1 & Traffic Prediction Result (Score + Breakdown) & High \\
2 & Time Slot Suggestions & Average \\
3 & Citizen Score Display & Low \\
4 & Admin Analytics Dashboard & High \\
5 & Complaint Status Report & Average \\
\bottomrule
\end{tabular}
\end{center}

\subsubsection*{External Inquiries (EQ)}
\begin{center}
\begin{tabular}{clc}
\toprule
\textbf{\#} & \textbf{Function} & \textbf{Complexity} \\
\midrule
1 & View Real-time Traffic Map & Average \\
2 & Search Route Information & Average \\
3 & View Complaint History & Low \\
4 & View Pothole Reports & Low \\
5 & View User List (Admin) & Low \\
\bottomrule
\end{tabular}
\end{center}

\subsubsection*{Internal Logical Files (ILF)}
\begin{center}
\begin{tabular}{clc}
\toprule
\textbf{\#} & \textbf{Data Store} & \textbf{Complexity} \\
\midrule
1 & Users Collection & Low \\
2 & PathInfo (Traffic Scores) & High \\
3 & Complaints Collection & Average \\
4 & Images (Pothole Reports) & Average \\
5 & Events Collection & Low \\
6 & Construction Collection & Low \\
7 & Diversions Collection & Low \\
8 & HotspotLocations Collection & Average \\
9 & PredictionLogs Collection & Average \\
10 & TrafficStatus Collection & Average \\
\bottomrule
\end{tabular}
\end{center}

\subsubsection*{External Interface Files (EIF)}
\begin{center}
\begin{tabular}{clc}
\toprule
\textbf{\#} & \textbf{External System} & \textbf{Complexity} \\
\midrule
1 & Google Maps Directions API & High \\
2 & OpenWeatherMap API & Average \\
3 & Calendarific Events API & Low \\
4 & MobileNet ML Model & High \\
5 & Cloudinary Image Storage & Average \\
\bottomrule
\end{tabular}
\end{center}

\subsection{Step 2: Calculate Unadjusted Function Points (UFP)}

The weighting factors for each complexity level are:

\begin{center}
\begin{tabular}{lccc}
\toprule
\textbf{Type} & \textbf{Low} & \textbf{Average} & \textbf{High} \\
\midrule
External Inputs (EI) & 3 & 4 & 6 \\
External Outputs (EO) & 4 & 5 & 7 \\
External Inquiries (EQ) & 3 & 4 & 6 \\
Internal Logical Files (ILF) & 7 & 10 & 15 \\
External Interface Files (EIF) & 5 & 7 & 10 \\
\bottomrule
\end{tabular}
\end{center}

\bigskip
\noindent\textbf{Calculation:}

\begin{center}
\begin{tabular}{l c c c c}
\toprule
\textbf{Type} & \textbf{Low} & \textbf{Average} & \textbf{High} & \textbf{Total} \\
\midrule
EI & $4 \times 3 = 12$ & $4 \times 4 = 16$ & $1 \times 6 = 6$ & \textbf{34} \\
EO & $1 \times 4 = 4$ & $2 \times 5 = 10$ & $2 \times 7 = 14$ & \textbf{28} \\
EQ & $3 \times 3 = 9$ & $2 \times 4 = 8$ & --- & \textbf{17} \\
ILF & $4 \times 7 = 28$ & $5 \times 10 = 50$ & $1 \times 15 = 15$ & \textbf{93} \\
EIF & $1 \times 5 = 5$ & $2 \times 7 = 14$ & $2 \times 10 = 20$ & \textbf{39} \\
\midrule
\multicolumn{4}{r}{\textbf{Unadjusted Function Points (UFP)}} & $\mathbf{211}$ \\
\bottomrule
\end{tabular}
\end{center}

\subsection{Step 3: Value Adjustment Factor (VAF)}

The VAF is calculated using 14 General System Characteristics (GSCs), each rated on a scale of 0--5.

\begin{center}
\begin{tabular}{clc}
\toprule
\textbf{\#} & \textbf{Characteristic} & \textbf{Rating (0--5)} \\
\midrule
1 & Data Communications & 5 \\
2 & Distributed Data Processing & 3 \\
3 & Performance & 4 \\
4 & Heavily Used Configuration & 2 \\
5 & Transaction Rate & 3 \\
6 & Online Data Entry & 5 \\
7 & End-User Efficiency & 4 \\
8 & Online Update & 4 \\
9 & Complex Processing & 4 \\
10 & Reusability & 3 \\
11 & Installation Ease & 3 \\
12 & Operational Ease & 3 \\
13 & Multiple Sites & 4 \\
14 & Facilitate Change & 3 \\
\midrule
 & \textbf{Total Degree of Influence (TDI)} & \textbf{50} \\
\bottomrule
\end{tabular}
\end{center}

\[
\text{VAF} = 0.65 + (0.01 \times \text{TDI}) = 0.65 + (0.01 \times 50) = \mathbf{1.15}
\]

\subsection{Step 4: Adjusted Function Points (AFP)}

\[
\text{AFP} = \text{UFP} \times \text{VAF} = 211 \times 1.15 = \mathbf{242.65 \approx 243~FP}
\]

\subsection{Step 5: Effort and Cost Estimation}

Using industry averages for a web application project:
\begin{itemize}[noitemsep]
    \item Productivity rate: \textbf{10 FP per person-month} (for a MERN stack project of moderate complexity)
    \item Average developer cost: \textbf{\textrm{INR} 40,000 per month} (academic/intern-level team)
\end{itemize}

\[
\text{Effort} = \frac{\text{AFP}}{\text{Productivity}} = \frac{243}{10} = \mathbf{24.3~\text{person-months}}
\]

With a team of 4 developers:
\[
\text{Duration} = \frac{24.3}{4} \approx \mathbf{6.1~\text{months}}
\]

\[
\text{Total Cost} = 24.3 \times 40{,}000 = \mathbf{\textrm{INR}~9,72,000}
\]

% ═══════════════════════════════════════════════════════════
% PART B: USE CASE POINTS
% ═══════════════════════════════════════════════════════════
\section{Use Case Points (UCP) Estimation}

\subsection{Overview}
The Use Case Points method, proposed by Gustav Karner in 1993, estimates effort based on the complexity of actors and use cases, adjusted by technical and environmental factors.

\subsection{Step 1: Unadjusted Actor Weight (UAW)}

Actors are classified by their interaction type:
\begin{center}
\begin{tabular}{llcc}
\toprule
\textbf{Actor} & \textbf{Type} & \textbf{Weight} & \textbf{Count} \\
\midrule
Citizen User & Complex (GUI-based) & 3 & 1 \\
Municipal Admin & Complex (GUI-based) & 3 & 1 \\
Google Maps API & Simple (API) & 1 & 1 \\
OpenWeatherMap API & Simple (API) & 1 & 1 \\
Calendarific API & Simple (API) & 1 & 1 \\
MobileNet ML Model & Average (protocol-driven) & 2 & 1 \\
Cloudinary & Simple (API) & 1 & 1 \\
\midrule
\multicolumn{3}{r}{\textbf{UAW}} & $\mathbf{12}$ \\
\bottomrule
\end{tabular}
\end{center}

\subsection{Step 2: Unadjusted Use Case Weight (UUCW)}

Use cases are classified by transaction count:
\begin{center}
\begin{tabular}{llcc}
\toprule
\textbf{Use Case} & \textbf{Complexity} & \textbf{Weight} & \textbf{Transactions} \\
\midrule
Register Account & Simple ($\leq$3) & 5 & 3 \\
Login / Logout & Simple & 5 & 2 \\
View Traffic Prediction & Complex ($>$7) & 15 & 12 \\
Get Time Suggestions & Average (4--7) & 10 & 5 \\
Upload Pothole Image & Complex & 15 & 8 \\
Submit Complaint & Average & 10 & 5 \\
Register Event Space & Simple & 5 & 3 \\
View Citizen Score & Simple & 5 & 2 \\
View Traffic Map & Average & 10 & 6 \\
Add Construction Data & Simple & 5 & 3 \\
Add Traffic Diversion & Simple & 5 & 3 \\
Review Citizen Reports & Average & 10 & 5 \\
Manage Users & Average & 10 & 6 \\
View Analytics Dashboard & Complex & 15 & 10 \\
\midrule
\multicolumn{2}{r}{\textbf{UUCW}} & $\mathbf{125}$ & \\
\bottomrule
\end{tabular}
\end{center}

\subsection{Step 3: Unadjusted Use Case Points (UUCP)}

\[
\text{UUCP} = \text{UAW} + \text{UUCW} = 12 + 125 = \mathbf{137}
\]

\subsection{Step 4: Technical Complexity Factor (TCF)}

\begin{center}
\begin{tabular}{clcc}
\toprule
\textbf{T\#} & \textbf{Factor} & \textbf{Wt.} & \textbf{Rating (0--5)} \\
\midrule
T1 & Distributed System & 2 & 4 \\
T2 & Response Time / Performance & 1 & 4 \\
T3 & End-User Efficiency & 1 & 4 \\
T4 & Complex Internal Processing & 1 & 5 \\
T5 & Reusability & 1 & 3 \\
T6 & Easy to Install & 0.5 & 3 \\
T7 & Easy to Use & 0.5 & 4 \\
T8 & Portability & 2 & 4 \\
T9 & Easy to Change & 1 & 3 \\
T10 & Concurrency & 1 & 3 \\
T11 & Security Features & 1 & 4 \\
T12 & Third-Party Access & 1 & 5 \\
T13 & Training Required & 1 & 2 \\
\midrule
\multicolumn{3}{r}{\textbf{TFactor} $= \sum(\text{Wt.} \times \text{Rating})$} & \textbf{52.5} \\
\bottomrule
\end{tabular}
\end{center}

\[
\text{TCF} = 0.6 + (0.01 \times 52.5) = \mathbf{1.125}
\]

\subsection{Step 5: Environmental Complexity Factor (ECF)}

\begin{center}
\begin{tabular}{clcc}
\toprule
\textbf{E\#} & \textbf{Factor} & \textbf{Wt.} & \textbf{Rating (0--5)} \\
\midrule
E1 & Familiarity with Project & 1.5 & 3 \\
E2 & Application Experience & 0.5 & 3 \\
E3 & OOP Experience & 1 & 4 \\
E4 & Lead Analyst Capability & 0.5 & 3 \\
E5 & Motivation & 1 & 4 \\
E6 & Stable Requirements & 2 & 3 \\
E7 & Part-Time Workers & $-1$ & 3 \\
E8 & Difficult Programming Language & $-1$ & 1 \\
\midrule
\multicolumn{3}{r}{\textbf{EFactor} $= \sum(\text{Wt.} \times \text{Rating})$} & \textbf{17.5} \\
\bottomrule
\end{tabular}
\end{center}

\[
\text{ECF} = 1.4 + (-0.03 \times 17.5) = 1.4 - 0.525 = \mathbf{0.875}
\]

\subsection{Step 6: Adjusted Use Case Points (UCP)}

\[
\text{UCP} = \text{UUCP} \times \text{TCF} \times \text{ECF} = 137 \times 1.125 \times 0.875 = \mathbf{134.86 \approx 135~UCP}
\]

\subsection{Step 7: Effort and Cost Estimation}

Using the standard productivity factor of \textbf{20 person-hours per UCP}:

\[
\text{Effort (hours)} = 134.86 \times 20 = \mathbf{2{,}697.2~\text{person-hours}}
\]

\noindent Converting to person-months (assuming 160 working hours/month):
\[
\text{Effort (months)} = \frac{2697.2}{160} = \mathbf{16.86~\text{person-months}}
\]

With a team of 4 developers:
\[
\text{Duration} = \frac{16.86}{4} \approx \mathbf{4.2~\text{months}}
\]

\[
\text{Total Cost} = 16.86 \times 40{,}000 = \mathbf{\textrm{INR}~6,74,400}
\]

% ═══════════════════════════════════════════════════════════
\section{Comparison of Results}
% ═══════════════════════════════════════════════════════════

\begin{center}
\begin{tabular}{lcc}
\toprule
\textbf{Parameter} & \textbf{FPA} & \textbf{UCP} \\
\midrule
Software Size & 243 FP & 135 UCP \\
Effort & 24.3 person-months & 16.9 person-months \\
Duration (4 developers) & 6.1 months & 4.2 months \\
Estimated Cost & INR 9,72,000 & INR 6,76,000 \\
\bottomrule
\end{tabular}
\end{center}

\medskip
\noindent The FPA method yields a higher estimate because it accounts for all internal data stores (10 MongoDB collections) and external interface files (5 APIs), which contribute significantly to the function point count. The UCP method, being use-case driven, focuses on user-facing interactions and tends to produce a more conservative estimate for backend-heavy systems.

\noindent The realistic estimate for this project lies between the two values, approximately \textbf{20 person-months} of total effort, translating to roughly \textbf{5 months} of development with a 4-member team.

% ═══════════════════════════════════════════════════════════
\section{Conclusion}
% ═══════════════════════════════════════════════════════════
This assignment demonstrated two complementary approaches to software cost estimation for the Traffic Prediction and Analysis System. Function Point Analysis provided a comprehensive size measurement accounting for all data interactions and system boundaries, yielding 243 adjusted function points. The Use Case Points method offered an effort-centric estimate grounded in use case complexity and team environment factors, producing 129 adjusted use case points.

Together, these techniques provide upper and lower bounds for project planning, enabling the development team to set realistic milestones, allocate resources proportionally, and manage stakeholder expectations with data-driven estimates.

\end{document}
